\documentclass{article}
\usepackage{kotex}
\usepackage{setspace}  % 줄간격
\setstretch{1.15}
\usepackage{enumitem}  % 목록 들여쓰기 
\setlist[itemize]{leftmargin=*, align=left}
\setlist[description]{
  font=\normalfont,
  labelwidth=4cm,
  leftmargin=4cm,
  labelsep=0cm,
  align=left
}
\usepackage{titlesec}  % 섹션 스타일 조정
\titleformat{\section}
  {\normalfont\Large\bfseries}
  {}
  {0pt}
  {}
  [\vspace{0.3em}\titlerule]

\titlespacing*{\section}
  {0pt}
  {1.5em}
  {0.8em}
\newcommand{\descriptionrule}{%
  \item[]\vspace{-0.6em}
  \noindent\hspace*{-\leftmargin}%
  \rule{0.6\linewidth}{0.4pt}
  \vspace{-0.4em}
}              
  \usepackage[a4paper, margin=2.5cm]{geometry}
\usepackage{fancyhdr}                 
\pagestyle{fancy}
\fancyhf{}                             
\fancyfoot[L]{\nouppercase{페미니즘 미학과 예술 (2026학년도 2학기)}} 
\fancyfoot[R]{\thepage}                
\renewcommand{\headrulewidth}{0pt}    
\renewcommand{\footrulewidth}{0.4pt}   

\begin{document}

\title{페미니즘 미학과 예술 (2026학년도 2학기)}
\author{\textit{written by using} \LaTeX}
\date{\today}
\maketitle
\thispagestyle{fancy}
\setcounter{secnumdepth}{0}
\bigskip

\section{0. 강좌정보}
\medskip
교과목번호: C20.144(001)\\
개설학과: 서울대학교 인문대학 미학과\\
담당교수: 안소현\\
수업시간: 화/목 09:30--10:45\\
수업장소: 1-103

\bigskip
\section{1. 수업목표}
\medskip
페미니즘 미학이 질문을 통해 어떻게 미학과 예술을 확장해 왔으며, 나아가 우리로 하여금 미래를 상상하게 하는지를 현대 미술을 중심으로 살펴본다.

\bigskip
\section{2. 교재 및 참고문헌}
\medskip
\subsection{교재}
『페미니즘 미학 입문 = Feminism』(Carolyn Korsmeyer, 2009.)\\
『왜 위대한 여성 미술가는 없었는가?』(Linda Nochlin, 2021.)\\
『몸 페미니즘을 향해: 무한히 변화하는 몸』(E. A. Grosz, 2019.)\\
『내셔널리즘과 젠더 = Nationalism to gender』(Ueno Chizuko, 2000.)\\
『젠더 트러블: 페미니즘과 정체성의 전복』(Judith Butler, 2008.)\\
『에코페미니즘』(Maria Mies et al., 2020.)\\
『해러웨이 선언문: 인간과 동물과 사이보그에 관한 전복적 사유』(Donna Jeanne Haraway, 2019.)\\
『트러블과 함께하기: 자식이 아니라 친척을 만들자』(Donna Jeanne Haraway, 2021.)\\
『변신: 되기의 유물론을 향해』(Rosi Braidotti, 2020.)\\
『오늘의 SF』(김원영 외., 2019.)

\newpage
\section{3. 강의계획}
\medskip
\begin{description}
    \item[강의개요] `아름다움', `취미', `예술', `천재' 등 근대 미학의 주요 개념들은 의외로 누구에게나 고르게 분배된 것이 아니었다. 다른 말로 하면 그것들은 이데올로기적이었고, 때로는 매우 성별적이었다. 따라서 이런 미학적 개념에 질문을 던지는 이들이 나타났고, 그 질문들은 새로운 예술을 새롭게 볼 수 있는 자리를 마련해주었다. 이 강의에서는 미학의 기본 개념들이 감춘 성별적 이데올로기들에 대한 질문에서 출발하여, 페미니즘 미학이 열어놓은 새로움을 현대 미술을 중심으로 살펴보고, 나아가 페미니스트들의 탈인간중심적 사유가 상상하게 하는 미래를 함께 엿볼 것이다.
    \item[강의방법] 페미니즘 미학의 중요한 텍스트를 참고할 뿐만 아니라, 다양한 현대 미술 작품들을 함께 살펴보고 비평적으로 토론하는 수업을 진행한다. 작품 이미지, 비디오 영상, 퍼포먼스 기록 등 가능한 많은 시각자료를 제시하고 적극적 토론을 유도함으로써 수강생이 페미니즘 미학의 스펙트럼을 최대한 폭넓은 관점에서 접하고 생각할 수 있도록 한다.
\end{description}
\medskip
\subsection{주차별 강의계획}
\medskip
\begin{description}
    \item[1회차(9/1)] 강의 개요 및 오리엔테이션: 왜 굳이, 페미니즘 미학인가라는 물음
    \item[2회차(9/3)] 페미니즘 미학과 맥락 읽기 : 페미니즘 미학에서 맥락 읽기의 중요성과 미술의 `맥락화'
    \item[3--4회차(9/8, 9/10)] 아름다움과 취미 : `아름다움'과 `취미', 그리고 `무관심적 관조' 같은 근대 미학의 핵심적 개념들이 어떤 이데올로기를 은폐해 왔는지 살펴본다.
    \item[5--6회차(9/15, 9/17)] 예술과 천재: `순수 예술(fine art)'이라는 개념이 무엇을 어떻게 배제해왔는가? 왜 위대한 여성 예술가는 없었는가? 예술가들을 신화화하면서 밀려나고 잃어버린 것들, 그에 대한 여성 예술가들의 대응을 따라가 본다.
    \item[7--8회차(9/22, 9/29)]  미술제도 : 예술은 개인적 감상의 대상에 그치지 않고, 제도의 자기장 안에서 규정되고 작동한다. 미술관, 비평 등 미술제도의 고르지 못함과 여성 미술가들의 문제제기에 주목한다.\footnote{9/24(추석연휴) 휴강}
    \item[9--10회차(10/1, 10/6)] 몸, 섹슈얼리티, 이미지 : 몸은, 특히 여성의 몸은 왜 문제적인가? 예술은 왜 그런 몸을 끊임없이 재현하려 하는가? 예술가들은 왜 섹슈얼리티에 주목하며, 자신의 몸을 드러냄으로써 무엇을 말하고자 하는가?
    \item[11회차(10/8)] 토론 : 섹슈얼리티와 공공성 (발리 엑스포트)
    \item[12--13회차(10/13, 10/15)] 숭고와 역겨운 것(abject) : 현대 미술은 왜 굳이 불편함을 자처하는가? 불편한 것들이 던지는 미학적 질문과 관련 여성주의 담론을 짚어본다.
    \item[14회차(10/20)] 새로운 매체, 새로운 재료, 새로운 이야기 : 매체, 재료, 기법도 성별 정치적이 될 수 있다. 음식, 공예, 퀼트, 여성 노동 등 평가절하된 것들을 무대에 올리려는 예술가들을 만난다.
    \item[15회차(10/22)] \textbf{중간고사}
    \newpage
    \item[16--17회차(10/27, 10/29)] 언어와 상징 : 언어는 페미니스트들에게 매우 중요한 지표이자 변화를 위한 수단이었다. 여성적 글쓰기를 비롯해 언어를 통해 은폐된 것들을 드러내려 한 예술가들을 만나볼 것이다
    \item[18회차(11/3)] 토론: 폭력을 재현하는 예술 (김홍석)
    \item[19회차(11/5)] 현장학습: 전시 단체 관람
    \item[20--21회차(11/10, 11/12)] 탈근대, 탈식민 : 근대국가 비판은 어떻게 페미니즘과 만나게 되었는가? 인종주의, 식민주의, 국가주의, 가부장제 담론을 넘어서려는 페미니스트 미학자와 예술가들의 주요 주장들을 살펴본다.
    \item[22-23회차(11/17, 11/19)] 젠더 트러블 - 젠더 정체성은 만들어진 것이다 : 현대의 정체성 담론은 우리가 당연한 경계라고 생각해온 것들을 뒤흔들었다. 열린 정체성을 외치는 예술가들과 최근의 퀴어 예술에 관한 논의에 대해 토론한다.
    \item[24--25회차(11/24, 11/26)] 에코페미니즘 : 누가 자연을 우리의 적으로 만들었는가 : 페미니스트의 타자에 대한 열린 시선은 인간중심주의를 넘어서는 새로운 세계를 열어놓았다. 기후위기 시대를 마주한 페미니스트들의 참신한 선언과 대안을 읽는다.
    \item[26--27회차(12/1, 12/3)] 포스트휴먼 이후 : 유토피아와 디스토피아 사이의 `SF의 세계', 사이보그 등 새로운 `크리처'에 대한 상상이 제기하는 페미니스트 미학의 관점들과 예술적 실험을 짚어본다.
    \item[28회차(12/8)] 토론: 사회적 논란과 재현의 윤리 (정윤석)
    \item[29회차(12/10)] 신유물론 페미니즘: 물질에 대한 새로운 이해를 통해 신유물론자들이 전개하는 페미니즘적 시선을 살펴본다.
    \item[30회차(12/15)] \textbf{기말고사}   
\end{description}

\bigskip
\section{4. 평가방법}
\medskip
\begin{description}
    \item[\textbf{성적부여방식}] 상대평가(A--F)
    \item[\textbf{출석(10\%)}] 수업일수의 1/3을 초과하여 결석하면 성적은 ``F'' 또는 ``U''가 됨\\(담당교수가 불가피한 결석으로 인정하는 경우는 예외로 할 수 있음)
    \item[\textbf{중간(40\%)}] 지필 논술형
    \item[\textbf{기말(50\%)}] 지필 논술형
    \descriptionrule
    \item[\textbf{합계}] 100\%
\end{description}

\newpage
\section{5. 생성형 AI 도구 활용 방침}
\medskip
본 강의에서는 생성형 AI를 학습 보조 도구로 활용할 수 있습니다. 시험 및 토론 준비에 필요한 아이디어 및 주제 탐색, 자료 분석 등의 목적으로 생성형 AI를 사용할 수 있습니다. 다만, 생성형 AI를 ‘주제 선정’에 사용한 경우, 사용 도구명, 사용 과정 및 생성 내용 등을 명시해야 합니다.

\bigskip
\section{6. 정원 외 신청 수용 여부}
\medskip
최대 10명

\bigskip
\section{7. 장애학생 지원사항}
\medskip
\subsection{강의 수강 관련}
\medskip
\begin{itemize}
    \item 시각장애: 교재 제작(디지털교재, 점자교재, 확대교재 등), 대필도우미 허용
    \item 지체장애: 교재 제작(디지털교재), 대필도우미 및 수업보조 도우미 허용
    \item 청각장애: 대필 및 문자통역 도우미 활동 허용, 강의 녹취 허용
    \item 건강장애: 질병 등으로 인한 결석에 대한 출석 인정, 대필도우미 허용
    \item 학습장애: 대필도우미 허용
    \item 지적장애/자폐성장애: 대필도우미 및 수업 멘토 허용
\end{itemize}
\subsection{과제 및 평가 관련}
\medskip
\begin{itemize}
    \item 시각장애/지체장애/청각장애/건강장애/학습장애: 과제 제출기한 연장, 과제 제출 및 응답 방식의 조정, 평가 시간 연장, 평가 문항 제시 및 응답 방식의 조정, 별도 고사실 제공
    \item 지적장애/자폐성장애: 개별화 과제 제출 및 대체 평가 실시
\end{itemize}
\subsection{비고}
\medskip
본 강의를 수강하는 장애학생들에게는 이상의 지원 서비스 이외에도 장애학생 개개인의 특성과 요구에 따라, 지도교수 및 장애학생지원센터와의 상담을 통하여 적절한 수준의 지원 서비스를 제공합니다. 장애학생에 대한 지원서비스와 관련하여 문의사항이 있는 학생들은 담당교수 안소현 혹은 장애학생지원센터(02-880-8787)로 문의바랍니다.

\end{document}


